% ===================================================================
%  LENTELĖ Nr. 2 — Žmogaus kūnas. — Pertrauka. — Žaidimai.
%  Traduction 2026 d'apres texto/io, controlee sur texto/fr, texto/ru
%  et texto/cs. Consigne : voir texto/lt/10-lentele-01.tex.
%
%  LE TABLEAU DU CORPS EST LE BANC D'ESSAI DE CETTE COLONNE, et il
%  donne d'emblee la mesure que le galicien avait apprise a lire :
%  la distance a la langue voisine se mesure, et elle est maximale
%  au corps.
%
%  MAIS LA MESURE SE RENVERSE ICI, ET C'EST L'INTERET DE LA COLONNE.
%  Pour le galicien, le corps etait l'endroit ou la langue voisine
%  cessait de servir de modele. Pour le lituanien, le corps est
%  l'endroit ou il n'a JAMAIS EU BESOIN d'un modele : « akis »
%  l'oeil, « ausis » l'oreille, « dantis » la dent, « širdis » le
%  coeur, « kraujas » le sang, « kaulas » l'os se lisent presque
%  comme le latin oculus, auris, dens, cor, cruor et le vieux
%  slave — non parce que le lituanien les a empruntes, mais parce
%  qu'il les a GARDES quand les autres les changeaient.
%
%  DE SORTE QUE LE TABLEAU 2 NE COUTE PAS UN SEUL EMPRUNT. Trente-
%  huit noms de parties du corps, et pas un mot international : le
%  cerveau est « smegenys », le crane « kaukolė », les poumons
%  « plaučiai », l'estomac « skrandis », les intestins « žarnos », la
%  cheville « kulkšnis », le mollet « blauzda », le talon « kulnas ».
%  Seuls « muskulas », « vena », « arterija » entrent de l'exterieur,
%  et ce sont les trois mots que le fac-simile lui-meme donne comme
%  savants. Le tcheque, treizieme colonne, avait du DECIDER de
%  refuser l'allemand ; le lituanien n'a rien a refuser au tableau
%  du corps, parce qu'il n'a rien pris.
%
%  LES JEUX, EN REVANCHE, SONT L'AUTRE MOITIE DU TABLEAU, et la
%  proportion s'inverse d'un coup. La note du feuillet 17 le dit
%  elle-meme : les jeux d'enfants portent des noms « pure
%  nacionala », et l'ido s'autorise soit a nommer par l'action, soit
%  a prendre le nom national. Le lituanien fait les deux dans la
%  meme page : « slėpynės », « vytynės », « šokdynė », « vilkelis »,
%  « kojūkai », « aitvaras », « sūpynės » sont a lui ; « kroketas »,
%  « tenisas », « futbolas », « boksuojasi », « fechtuojasi »,
%  « trapecija » viennent du dehors avec la chose. Le corps ne
%  s'importe pas, le sport si — et c'est le meme critere du lieu
%  d'apprentissage, applique non plus a un mot mais a un tableau
%  entier.
%
%  DEUX NOMS COMPOSES PAR L'ACTION, comme la note y invite. Le
%  « frap-divino » de l'ido, la main-chaude du francais, devient
%  « smūgio spėlionės », la devinette du coup ; et le couple
%  « shultro-salto » / « dorso-salto », que la note distingue par
%  l'endroit ou se posent les mains, devient « pečių šuolis » /
%  « nugaros šuolis », le saut des epaules et le saut du dos. Nommer
%  par l'action garde la distinction que le fac-simile explique ;
%  prendre un nom national l'aurait perdue.
%
%  ZERO INVERSION, et le genitif antepose n'a pas eu l'occasion de
%  coliger deux renvois : au tableau du corps les mots se suivent
%  sans se subordonner, et aux jeux le nom du jeu porte seul son
%  renvoi. Total apres deux tableaux : 0.
% ===================================================================

%%K t02-apar-1 apar
\VUsaut{4.0mm}
\VUtitre{15.0pt}{60}{LENTELĖ Nr. 2}
\VUsaut{2.0mm}
\VUtitre{11.4pt}{0}{Žmogaus kūnas. --- Pertrauka.}
\VUtitre{11.4pt}{0}{Žaidimai.}

%%K t02-tit-0 sub
\VUtitre{13.2pt}{0}{Žmogaus kūnas.}
\VUsaut{2.4mm}
\VUcentre{10.2pt}{0}{\textit{(Paprasta gamtos mokslo pamoka.)}}
\VUsaut{3.0mm}

%%K t02-01-1 p
1. --- Pagrindinės žmogaus kūno dalys yra: \VUgras{galva,} \VUgras{liemuo} ir
\VUgras{galūnės.}

%%K t02-tit-1 sub
\VUpk{10.2pt}{40}{Galva.}
\VUsaut{3.0mm}

%%K t02-02-1 p
2. --- Galvą sudaro dvi dalys: \VUgras{kaukolė,} kurioje yra \VUgras{smegenys} ir
kurią dengia \VUgras{plaukai,} ir \VUgras{veidas.}

%%K t02-03-1 p
3. --- Mūsų piešinyje nematyti nei kaukolės, nei smegenų; bet labai aiškiai matome
svarbiąsias veido dalis.

%%K t02-04-1 p
4. --- Štai pirmiausia \VUgras{kakta,} rodanti galingą protą, nuolat dirbančią
mintį. \VUgras{Smilkiniai} labai laisvi. \VUgras{Akis} gyva ir plačiai atmerkta. Ją
saugo tankus \VUgras{antakis} ir ilgos \VUgras{blakstienos.}

%%K t02-05-1 p
5. --- Štai dar \VUgras{nosis} ir \VUgras{šnervės.} Nosis mums padeda uosti ir
kvėpuoti. Abipus nosies yra \VUgras{skruostai,} o toliau \VUgras{ausys.} Apatinę
veido dalį sudaro \VUgras{burna} ir \VUgras{smakras.}

%%K t02-06-1 p
6. --- Jų ne itin gerai matome, nes juos nuo mūsų slepia \VUgras{ūsai} ir
\VUgras{žandenos.} Bet žinome, kad burnoje yra dvi \VUgras{lūpos,} \VUgras{viršutinis}
ir \VUgras{apatinis žandikaulis} su savo \VUgras{dantimis} ir \VUgras{liežuvis.} Burna
kalbame ir valgome. Liežuviu jaučiame maisto skonį.

%%K t02-07-1 p
7. --- Akis, ausis, nosis ir burna, kaip ir pirštai, yra penkių juslių organai:
\VUgras{rega,} \VUgras{klausa,} \VUgras{uoslė,} \VUgras{skonis,} \VUgras{lytėjimas.}

%%K t02-08-1 p
8. --- Taigi galva yra viena įdomiausių žmogaus kūno dalių.

%%K t02-tit-2 sub
\VUsaut{4.4mm}
\VUpk{10.2pt}{40}{Liemuo.}
\VUsaut{3.0mm}

%%K t02-09-1 p
9. --- \VUgras{Kaklas} jungia galvą su liemeniu. Jį sudaro \VUgras{gerklė} ir
\VUgras{sprandas.}

%%K t02-10-1 p
10. --- Liemenyje taip pat yra daug svarbių organų. \VUgras{Krūtinėje} yra
\VUgras{plaučiai,} kuriais kvėpuojame, ir \VUgras{širdis,} kuri yra gyvybės centras.
Kai širdis nustoja plakti, gyva būtybė miršta.

%%K t02-11-1 p
11. --- \VUgras{Šonkauliai} laiko krūtinę.

%%K t02-12-1 p
12. --- Kitos liemens dalys yra \VUgras{šonas,} \VUgras{nugara,} \VUgras{pečiai} ir
\VUgras{pilvas.} Miegoti gulame ant šono arba ant nugaros; naštas nešame ant peties.

%%K t02-13-1 p
13. --- Pilve yra \VUgras{skrandis} ir \VUgras{žarnos.} Jie padeda mums virškinti
maistą.

%%K t02-tit-3 sub
\VUsaut{4.4mm}
\VUpk{10.2pt}{40}{Galūnės.}
\VUsaut{3.0mm}

%%K t02-14-1 p
14. --- Turime keturias \VUgras{galūnes}: dvi \VUgras{rankas} ir dvi \VUgras{kojas.}
Rankomis imame, laikome, keliame, renkame daiktus; kojomis stovime, einame ir bėgame.

%%K t02-15-1 p
15. --- \VUgras{Petys} jungia ranką su kūnu. Labai tvirtas \VUgras{raumuo,}
dvigalvis, judina ranką, kurią \VUgras{alkūnė} jungia su \VUgras{dilbiu.}
\VUgras{Riešas} sudaro \VUgras{plaštakos} sąnarį, o ji baigiasi \VUgras{pirštais} ir
\VUgras{nagais.} Ant plaštakos įžiūrime \VUgras{veną} (ir arteriją), kuria teka
\VUgras{kraujas.}

%%K t02-16-1 p
16. --- Kojos dalys yra: \VUgras{šlaunis,} \VUgras{kelis} ir \VUgras{pakinklis,}
\VUgras{blauzda,} kurios \VUgras{mėsą} dengia \VUgras{oda,} \VUgras{kulkšnis} ir
\VUgras{pėda.} Kojos \VUgras{kaulai} labai stori ir labai stiprūs. Pėdos gale yra
\VUgras{kojų pirštai.} Kitos dalys yra \VUgras{padas} ir \VUgras{kulnas.}

%%K t02-17-1 p
17. --- Toks yra beveik išsamus žmogaus kūno aprašymas.

%%K t02-17-2 p
\textit{Pastaba.} --- \textit{Pasitelkdamas posakį: „man skauda... (galvą ir
t. t.)“, mokytojas galės dar kartą panaudoti visą šį žodyną, kalbėdamas apie gerą ar
blogą} \VUgras{kūno būklę.}

%%K t02-tit-4 sub
\VUsaut{5.0mm}
\VUpk{10.2pt}{40}{Gimnastika.}
\VUsaut{2.4mm}
\VUcentre{10.2pt}{0}{\textit{(Vieno mokinio pasakojimas.)}}
\VUsaut{3.0mm}

%%K t02-18-1 p
18. --- Kad būtume lankstūs, vikrūs ir sveiki, turime mankštinti kojas ir rankas.
Todėl žaidžiame ir mankštinamės \VUgras{gimnastika.}

%%K t02-19-1 p
19. --- Per savaitę abu šie pratimai vyksta licėjaus kieme (žiūrėk lentelę
Nr.~1). Bet ketvirtadieniais ir sekmadieniais einame į didelę pievą, kur turime
vietos bėgioti ir linksmintis.

%%K t02-20-1 p
20. --- Mūsų mokytojai ir tėvai kartais mus ten palydi; bet pratimams vadovauja
\VUgras{gimnastikos mokytojas} \textsuperscript{(12)}.

%%K t02-21-1 p
21. --- Pievoje, kairėje nuo įėjimo, yra \VUgras{gimnastikos įrenginiai}
\textsuperscript{(1)}: lipame \VUgras{kopėčiomis} \textsuperscript{(2)}, aukštyn
kojomis mituojame ant \VUgras{žiedų} \textsuperscript{(3)} ir ant \VUgras{trapecijos}
\textsuperscript{(4)}, rankomis keliamės iki \VUgras{virvinių kopėčių}
\textsuperscript{(5)} arba \VUgras{mazgotos virvės} \textsuperscript{(6)} viršaus.

%%K t02-22-1 p
22. --- Tuo tarpu mūsų draugai šokinėja per \VUgras{medinį arklį}
\textsuperscript{(7)} arba per \VUgras{lygiagretes} \textsuperscript{(11)}. Kiti
\VUgras{boksuojasi} \textsuperscript{(8)} arba \VUgras{fechtuojasi}
\textsuperscript{(9)} su kauke ir \VUgras{rapyra.} Dar kiti \VUgras{kilnoja
svarmenis} \textsuperscript{(10)}.

%%K t02-tit-5 sub
\VUsaut{4.6mm}
\VUpk{10.2pt}{40}{Žaidimai.}
\VUsaut{3.0mm}

%%K t02-23-1 p
23. --- Aplink gimnastikuojančius berniukai ir mergaitės žaidžia įvairius
\VUgras{žaidimus.} Mergaitės linksminasi su savo \VUgras{lanku}
\textsuperscript{(14)}; jos šokinėja per \VUgras{šokdynę} \textsuperscript{(13)} arba
kviečia savo brolius ir pusbrolius į \VUgras{kroketo} \textsuperscript{(15)} partiją.
Joms miela mediniu \VUgras{plaktuku} nustumti priešininko \VUgras{rutulį} kaip tik
tada, kai jis mano lengvai pravažiuosiąs pro lankelį!

%%K t02-24-1 p
24. --- Berniukai mieliau renkasi raumeningesnius žaidimus. Jie noriai žaidžia
\VUgras{kėgliais} \textsuperscript{(16)}, kuriuos vikriai nuverčia \VUgras{rutuliu}
\textsuperscript{(17)}. Jie taip pat nepaniekina \VUgras{vilkelio}
\textsuperscript{(18)} ir \VUgras{bilbokės} \textsuperscript{(19)} ar net
\VUgras{varlės žaidimo} \textsuperscript{(20)}. Bet labiausiai mėgstami jų žaidimai
yra \VUgras{rutuliukai} \textsuperscript{(21)} ir \VUgras{sieninis kamuoliukas}
\textsuperscript{(22)}, nes \VUgras{šiltnamio} \textsuperscript{(26)} siena labai
tinka lengvam kamuoliukui atšokdinti.

%%K t02-25-1 p
25. --- Mažasis Jonas, vaikščiodamas savo \VUgras{kojūkais} \textsuperscript{(24)},
labai vikrus ir puikiai moka išlaikyti pusiausvyrą. Netoli jo keli draugai žaidžia
\VUgras{slėpynių} \textsuperscript{(25)} tame šiltnamyje ir už \VUgras{gyvatvorės.}
Kiti mieliau renkasi \VUgras{smūgio spėliones} \textsuperscript{(28)} arba
\VUgras{skrajuką} \textsuperscript{(29)}, kuris lengvai plazda nuo vienos
\VUgras{raketės} prie kitos.

%%K t02-noto-1 noto
\VUnotes{143.2mm}{%
(*) Vaikų žaidimai dažnai turi grynai tautinius pavadinimus. Stengėmės kuo geriau
pavadinti juos ido kalba, arba semdamiesi įkvėpimo iš paties veiksmo, kuris juos
apibūdina, arba pasirinkdami tautinį vardą tame ar kitame krašte, kai jis nėra pernelyg
neteisingas. Pvz.: \VUgras{statinės žaidimas} (vokiečių ir prancūzų) yra
\VUgras{varlės žaidimas} \textsuperscript{(20)} (anglų), kuriame žaidėjai stengiasi
įmesti savo apvalius diskelius pro metalinės varlės, žiopsančios ant skylėtos dėžės
(statinės), burną.
\VUgras{Pečių šuolis} \textsuperscript{(30)} nuo \VUgras{nugaros šuolio} skiriasi tik
šiuo: norėdami šokti per galvą, žaidėjai remia rankas į savo partnerio pečius, o
„\VUgras{nugaros šuolyje}“ jie remia rankas tik į jo nugarą. --- Arba dar keli
dalyviai palinksta, o kiti šoks per jų nugaras, remdami rankas į petį; pirmieji turi
iškęsti smūgį nepalūždami, antrieji turi šokti neprarasdami pusiausvyros (žiūrėk pačią
lentelę).%
}

%%K t02-26-1 p
26. --- Pievos viduryje yra du medžiai. Prie vieno iš jų kojos septyni ar aštuoni
berniukai surengė \VUgras{pečių šuolio} \textsuperscript{(30)} (*) partiją. Ne
visuose kraštuose šis žaidimas žinomas; bet visi žino, kas yra \VUgras{nugaros
šuolis,} kuriuo vaikai šįkart nežaidžia.

%%K t02-tit-6 sub
\VUsaut{5.0mm}
\VUpk{10.2pt}{40}{Žaidimai gryname ore.}
\VUsaut{3.4mm}

%%K t02-27-1 p
27. \VUgras{Akluko žaidimas} \textsuperscript{(31)} taip pat labai linksmas;
mergaitės ir berniukai paeiliui užsiriša ant akių \VUgras{raištį} ir gaudo
nevikruosius, kurie leidžiasi pagaunami.

%%K t02-28-1 p
28. --- Pievos viduryje štai labai prancūziškas, bet kiek smarkus žaidimas: tai
\VUgras{meškos žaidimas} \textsuperscript{(32)}, daug triukšmingesnis už tokį ramų
\VUgras{aitvaro} \textsuperscript{(33)} žaidimą ar net už anglišką žaidimą
\VUgras{tenisą} \textsuperscript{(34)}.

%%K t02-29-1 p
29. --- Šis paskutinis sportas yra vyresniųjų mokinių džiaugsmas. Patys mūsų
mokytojai jį noriai žaidžia. Jis reikalauja didelio vikrumo ir miklumo ir man labai
patinka. \VUgras{Sūpynių} \textsuperscript{(35)} žaidimus palieku mergaitėms.

%%K t02-30-1 p
30. --- Nemėgstu kito angliško žaidimo, kuris galbūt taptų pavojingas.

%%K t02-30-1b p
Turiu galvoje \VUgras{futbolą} \textsuperscript{(36)}, kuris džiugina mano vyresnįjį
brolį. Man labiau patinka mūsų senosios \VUgras{vytynės} \textsuperscript{(38)}, tokios
pat sveikos ir tikrai mažiau žiaurios.

%%K t02-tit-7 sub
\VUsaut{4.6mm}
\VUpk{10.2pt}{40}{Dviratis ir baliono žaidimas.}
\VUsaut{3.0mm}

%%K t02-31-1 p
31. --- Taip pat noriai apvažiuoju pievą \VUgras{dviračiu} \textsuperscript{(37)}.

%%K t02-32-1 p
32. --- Kitais metais mokysiuos \VUgras{jodinėjimo} \textsuperscript{(39)}. Koks
gražus arklys, kuris grakščiai žingsniuoja ar risnoja ir kuris šuoliais peršoka
kliūtis!

%%K t02-33-1 p
33. --- Vasarą mokomės \VUgras{plaukimo} \textsuperscript{(40)}. Su malonumu nardome
ir maudomės skaidriame ir vėsiame upės vandenyje. Kartais galiausiai paleidžiame
popierinį \VUgras{balioną} \textsuperscript{(41)}, kuris didingai kyla į orą, vos tik
prisipildo šilto oro.

%%K t02-34-1 p
34. --- Yra dar daug kitų žaidimų, bet nebaigčiau jų vardyti, ir iš šios santraukos
suprantama, kad ketvirtadieniais mūsų gražioje pievoje nelabai nuobodžiaujama.

%%K t02-34-2 p
N.-B. --- \textit{Mokytojas lengvai papildys šį skyrių, išsamiau aprašydamas kiekvieną
iš svarbiausių žaidimų.}
